\label{sec:future}
\label{app:future}

\begin{table*}[t]
\centering
% Define a new column type for X columns to make text left-aligned and allow hyphenation
\newcolumntype{L}{>{\RaggedRight\arraybackslash}X}

\begin{tabularx}{\textwidth}{@{}lL L@{}} % Use \textwidth for table width, l for first col, L (custom X) for others
\toprule
                          & \textbf{Human}                                                                         & \textbf{Self-Referential Agent}                                                                 \\ \midrule
Intelligent Module        & brain                                                                                  & LLM                                                                                             \\
Perceptual and Action Module    & body                                                                                   & code and tool                                                                                   \\
Self-Referential Feature  & Humans can train their brain and body to improve, thus becoming better                 & Self-referential agents can modify their code, even the underlying LLM, to improve themselves \\
Self-Awareness Question           & Can the brain recognize itself as a brain? Can it perceive its own mode?               & Can LLM understand that it is one part of the modified codes?                               \\ \bottomrule
\end{tabularx}
\caption{An analogy of self-reference for both humans and agents}
\label{tab:self_reference_analogy}
\end{table*}


\subsection{Discussions}
%The distinction between learning and meta-learning is inherently ambiguous. Suppose a model is divided into a learning algorithm and a core model component. If the learning algorithm is modified through training, it can be viewed as meta-learning. However, if the entire model is considered as a unified entity, these modifications represent standard learning processes that adjust the model's parameters. This indistinguishability between learning and meta-learning suggests that artificially imposing hierarchical structures in meta-learning is unnatural. Instead, meta-learning optimization should be treated as a form of standard learning.

% Reflecting on the evolution of artificial intelligence, from expert systems to support vector machines, to various neural networks, and finally to scaling laws, it becomes evident that complex, human-crafted designs have often become bottlenecks in the development of AI. From diverse and intricate tasks to various pre-training tasks and next-word prediction, these shifts highlight a key insight: complex manual designs tend to limit AI's potential, whereas simpler, more generalizable designs have proven to unlock greater possibilities and scalability. \yxj{too abstract}
% \framework embodies this very philosophy.

Table~\ref{tab:self_reference_analogy} draws an analogy between human self-reference and the potential for self-referential capabilities in artificial agents. 
Inspired by this analogy, we believe that self-reference constitutes a foundational and indispensable attribute for the development of AGI, and that future agents should inherently be self-referential.
As foundation models grow in power, agents can more effectively enhance their own capabilities, ultimately evolving beyond the boundaries (or limitations) of human design.

Furthermore, when an agent adjusts its own code based on feedback, this is akin to an \textit{executable} version of test-time computing. In the context of LLMs, test-time computing typically involves generating additional tokens during inference, which then serve as a prefix to the final answer. This is because LLMs process information solely through text, making this their primary method for increasing computational effort at test time. For agents, however, their ability to call tools and execute code allows for far more diverse forms of test-time computing. \framework{} actualizes these more diverse forms of test-time computing precisely by modifying its own runtime code during test time.


\subsection{Future Directions}




There is significant room for improvement in the effectiveness, efficiency, and robustness of the \framework{}'s self-improvement capabilities, which requires better initial designs. The following are some promising directions for enhancement:
1) \textbf{Enhanced Optimization Modules}: Utilize human priors to design more effective optimization modules, such as genetic algorithms and reinforcement learning frameworks.
2) \textbf{Expanded Modifiability}: Broaden the scope of permissible modifications, allowing the agent to design and execute code that can fine-tune its own LLM modules.
3) \textbf{Improved Environmental Feedback and Task Sequencing}: Implement more sophisticated environmental feedback mechanisms and carefully curated task sequences during the initial optimization phase to prime the agent's capabilities. Once the agent demonstrates sufficient competence, it can then be exposed to real-world environments.



In addition, there are several other directions worth exploring and analyzing:

\noindent\textbf{Collective Intelligence} \quad Investigate the interactions among multiple \godel Agents. Agents could consider other agents as part of their environment, modeling them using techniques such as game theory. This approach treats these agents as predictable components of the environment, enabling the study of properties related to this specific subset of the environment. 

\noindent\textbf{Agent and LLM Characteristics} \quad Use the \framework{}’s self-improvement process as a means to study the characteristics of agents or LLMs. For example, can an agent genuinely become aware of its own existence, or does it merely analyze and improve its state as an external observer? 
This line of inquiry could yield insights into the nature of self-awareness in artificial systems.

\noindent\textbf{Theoretical Analysis} \quad Explore whether \framework{} can achieve theoretical optimality and what the upper bound of its optimization might be. Determine whether the optimization process could surpass the agent’s own understanding, and if so, at what point this might occur.

\noindent\textbf{Safety Considerations} \quad Although the current behavior of FMs remains controllable, as their capabilities grow, fully self-modifying agents will require human oversight and regulation. It may become necessary to limit the scope and extent of an agent's self-modifications, ensuring that modifications occur only within a controlled environment.



% \subsection{Limitations}
% As the first self-referential agent, \framework{} has to construct all task-related code autonomously, which poses significant challenges. Consequently, this work does not compare directly with the most complex existing agent systems, such as OpenDevin~\citep{wang2024opendevinopenplatformai}, which have benefited from extensive manual engineering efforts. Currently, \framework{} is not sufficiently stable and may be prone to error accumulation, hindering its ability to continue self-optimization. 
%This makes it unrealistic to expect it to outperform systems that have taken researchers several months or even years to develop. The experiments presented in this paper are intended to demonstrate the feasibility of recursive self-improvement. 
% A more robust and advanced implementation of the \framework{} is anticipated, with numerous potential improvements outlined in Section \ref{sec:future}. 
% The experiments presented in this paper are intended to demonstrate the effectiveness and feasibility of recursive self-improvement. 
% \yxj{too honest}

%Additionally, as the agent system becomes increasingly complex through self-optimization, it may require exponentially more intelligence to understand itself. Consequently, a system capable of complete self-referential at the outset may lose this capability as it evolves~\citep{Yampolskiylimitation}. The exact point at which the agent can no longer comprehend and improve itself has not been thoroughly explored. Investigating this phenomenon, both experimentally and theoretically, could provide valuable insights into the limitations of recursive self-improvement.

